\documentclass[12pt,a4paper]{article}
\usepackage[utf8]{inputenc}
\usepackage[portuguese]{babel}
\usepackage{graphicx}
\usepackage{hyperref}
\usepackage{geometry}
\geometry{a4paper, margin=2.5cm}

\title{\textbf{Documentação Oficial - Projeto Cinema APO2}}
\author{Guilherme Barroso Costa e Arthur Ferreira Fernandes}
\date{\today}

\begin{document}

\maketitle

\section{Objetivo do Sistema}
O sistema desenvolvido visa automatizar o processo de reserva de ingressos de cinema. Oferece uma interface web moderna onde clientes podem se cadastrar, validar seu e-mail, gerenciar seus perfis e reservar ingressos para as sessões em cartaz. Possui também uma área restrita para o administrador gerenciar o catálogo de filmes.

\section{Equipe e Funções}
\begin{itemize}
    \item \textbf{Guilherme Barroso Costa} (GU3054101): Desenvolvedor Fullstack e Banco de Dados.
    \item \textbf{Arthur Ferreira Fernandes} (GU3055621): QA, Documentação e Arquitetura J2EE.
\end{itemize}

\section{Requisitos do Sistema}
\subsection{Funcionais}
\begin{itemize}
    \item O sistema deve permitir o cadastro de usuários e envio de e-mail de validação.
    \item O sistema deve autenticar os usuários e proteger rotas privadas.
    \item O sistema deve permitir que clientes reservem ingressos para as sessões.
    \item O administrador deve poder realizar um CRUD do catálogo de filmes.
\end{itemize}

\subsection{Não Funcionais}
\begin{itemize}
    \item O sistema deve ser desenvolvido em J2EE (Servlets e JSP).
    \item O banco de dados utilizado será o MySQL 8.
    \item As requisições entre Front-End e Back-End devem ser assíncronas (AJAX).
    \item Os endpoints devem retornar unicamente o formato JSON.
    \item A interface deve ser estilizada com Bootstrap 5.
\end{itemize}

\section{Diagramas}
\textit{Nota: Insira as imagens exportadas na pasta imagens/ e remova os comentários abaixo.}
% \begin{figure}[h]
%     \centering
%     \includegraphics[width=0.8\textwidth]{imagens/diagrama_classes.png}
%     \caption{Diagrama de Classes}
% \end{figure}

\section{Relatório de Testes de Endpoints}

Foram implementados endpoints RESTful (retornando JSON) utilizando Servlets em conformidade com o padrão arquitetural MVC, testados através de requisições AJAX (jQuery) e ferramentas de teste de API (Postman).

\subsection{Endpoint: Autenticação (Login)}
\begin{itemize}
    \item \textbf{URL:} \texttt{/api/login}
    \item \textbf{Método HTTP:} \texttt{POST}
    \item \textbf{Parâmetros (x-www-form-urlencoded):} \texttt{email} (String), \texttt{senha} (String)
    \item \textbf{Resposta Esperada (Sucesso):}
    \begin{verbatim}
    { "status": "success", "message": "Login realizado com sucesso.", 
      "redirect": "privada/dashboard.jsp" }
    \end{verbatim}
    \item \textbf{Resposta Esperada (Erro):} \texttt{\{ "status": "error", "message": "Credenciais inválidas." \}}
\end{itemize}

\subsection{Endpoint: Cadastro de Clientes (Registro)}
\begin{itemize}
    \item \textbf{URL:} \texttt{/api/registro}
    \item \textbf{Método HTTP:} \texttt{POST}
    \item \textbf{Parâmetros:} \texttt{nome}, \texttt{cpf}, \texttt{telefone}, \texttt{email}, \texttt{senha}
    \item \textbf{Resposta Esperada:} \texttt{\{ "status": "success", "message": "Cadastro realizado com sucesso!" \}}
\end{itemize}

\subsection{Endpoint: Listagem de Filmes (Cliente)}
\begin{itemize}
    \item \textbf{URL:} \texttt{/api/filmes}
    \item \textbf{Método HTTP:} \texttt{GET}
    \item \textbf{Regra de Negócio:} Retorna exclusivamente filmes que possuem sessões ativas programadas para datas futuras (\texttt{s.horario >= NOW()}).
    \item \textbf{Resposta Esperada:} Array de objetos contendo \texttt{id, nome, duracao, genero, classificacao, sinopse}.
\end{itemize}

\subsection{Endpoint: Gerenciamento de Filmes e Alocação Inicial (Admin)}
\begin{itemize}
    \item \textbf{URL:} \texttt{/api/admin/filmes}
    \item \textbf{Método HTTP (GET):} Retorna o acervo completo de filmes (independentemente de possuírem sessões).
    \item \textbf{Método HTTP (POST):} Aceita dados do filme (\texttt{titulo, duracao, genero, classificacao, sinopse}). Caso sejam enviados também \texttt{sala\_id, horario} e \texttt{valor\_ingresso}, realiza a criação automática e simultânea da Sessão no mesmo escopo transacional.
    \item \textbf{Método HTTP (PUT/DELETE):} Atualiza ou remove filmes pelo \texttt{id}.
\end{itemize}

\subsection{Endpoint: Gerenciamento de Salas (Admin)}
\begin{itemize}
    \item \textbf{URL:} \texttt{/api/admin/salas}
    \item \textbf{Método HTTP (GET):} Retorna as salas cadastradas com atributos de \texttt{capacidade} e flag de \texttt{disponivel}.
    \item \textbf{Método HTTP (POST/PUT/DELETE):} CRUD completo para inserção, bloqueio ou alteração de capacidade.
\end{itemize}

\subsection{Endpoint: Gerenciamento de Sessões (Admin)}
\begin{itemize}
    \item \textbf{URL:} \texttt{/api/admin/sessoes}
    \item \textbf{Método HTTP (GET):} Retorna as sessões com join agilizado trazendo \texttt{filmeTitulo} e identificação da sala.
    \item \textbf{Método HTTP (POST/PUT/DELETE):} Alocação e reconfiguração de horários e valores de ingressos.
\end{itemize}

\subsection{Endpoint: Limpeza e Manutenção de Salas (Admin)}
\begin{itemize}
    \item \textbf{URL:} \texttt{/api/admin/limpeza} e \texttt{/api/admin/manutencao}
    \item \textbf{Método HTTP (GET):} Retorna a listagem de salas e o estado atual (\texttt{AGUARDANDO}, \texttt{EM ANDAMENTO}, \texttt{CONCLUIDO}).
    \item \textbf{Método HTTP (POST):} Atualiza o status e, no caso de início de manutenção, ajusta automaticamente a flag \texttt{disponivel = false} para travar vendas na sala.
\end{itemize}

\subsection{Endpoint: Mapeamento de Poltronas e Venda de Ingressos}
\begin{itemize}
    \item \textbf{URL Poltronas (GET):} \texttt{/api/poltronas?salaId=X} (Retorna array de assentos com flag \texttt{disponivel}).
    \item \textbf{URL Ingressos (POST):} \texttt{/api/privada/ingressos} (Recebe JSON contendo \texttt{sessaoId} e array de \texttt{poltronas} selecionadas, registrando o vínculo ao cliente logado na sessão).
\end{itemize}

\section{Relatório de Testes de Funcionalidades}

Esta seção detalha os testes manuais e funcionais realizados diretamente na interface do usuário (Views).

\begin{table}[h]
\centering
\begin{tabular}{|p{3cm}|p{5cm}|p{4cm}|p{3cm}|}
\hline
\textbf{Caso de Uso} & \textbf{Passo a Passo} & \textbf{Resultado Esperado} & \textbf{Status/Erros} \\ \hline
1. Login Usuário Válido & 1. Acessar tela de Login.\newline 2. Inserir email teste@cinema.com\newline 3. Inserir senha teste123\newline 4. Clicar em Entrar & O sistema exibe um alerta verde e redireciona para a área privada dashboard.jsp. & OK. Nenhum erro. \\ \hline
2. Controle de Acesso & 1. Copiar a URL da área administrativa.\newline 2. Acessar diretamente sem fazer login. & O sistema intercepta pelo AuthFilter e redireciona para a página de Login com erro de acesso. & OK. Nenhum erro. \\ \hline
3. Cadastro de Filme (Admin) & 1. Logar como Administrador.\newline 2. Ir em Gerenciar Filmes.\newline 3. Preencher campos e Salvar. & O filme aparece na listagem AJAX imediatamente sem recarregar a página. & OK. Nenhum erro. \\ \hline
4. Reserva de Ingresso (User) & 1. Logar como Cliente.\newline 2. Acessar 'Ver Filmes'.\newline 3. Clicar em Comprar em um filme.\newline 4. Confirmar quantidade. & O sistema salva no banco a reserva atrelada ao Cliente logado. & OK. Nenhum erro. \\ \hline
\end{tabular}
\caption{Tabela de Testes Funcionais da Interface}
\end{table}


\end{document}
